\chapter{Grundlagen der Brandentstehung}\label{ch:grundlagenBrandentstehung}

Die Entstehung eines Brandes setzt das Vorhandensein von drei Komponenten voraus, die im richtigen Mengenverhältnis vorliegen müsssen: brennbares Material, Sauerstoff als Oxidationsmittel und eine Zündtemperatur, die erreicht werden muss, um die chemische Kettenreaktion in Gang zu setzen. Brennbares Material kann dabei in verschiednene Formen auftreten, wie z.B. Holz, Laub oder trockenen Pflanzenreste. Sauerstoff wirkt als Oxidationsmittel und liegt in unserer Atmosphäre in ausreichender Menge vor. Die Zündtemperatur variiert je nach Material und hängt von Faktoren wie Feuchtigkeitsgehalt und Dichte ab. Wenn diese drei Komponenten zusammentreffen, beginnt der Verbrennungsprozess, der als chemische Reaktion der Oxidation eines Brandstoffes durch ein Oxidationsmittel in Gegenwart einer Energiequelle definiert ist. Energiequellen können in diesem Zusammenhang beispielsweise mechanische Reibung, Elektrizität, chemische Reaktionen oder Solarenergie sein.

Ein Waldbrand durchläuft in der Regel mehrere Phasen, beginnend mit der Enstehungsphase, auch als Schwelbrandphase bezeichnet, in der das Feuer noch klein und auf einen begrenzten Raum beschränkt ist. In dieser Phase sind die Materialien oft nur am Glimmen oder es bilden sich erste kleine Flammen, die jedoch noch keine offene Flamme darstellen. Die Hitze- und Rauchentwicklung sind hier noch relativ gering, da die Temperatur zu Beginn noch erträglich ist, allerdings ist der Rauch hochgiftig und sollte in keinem Fall unterschätzt werden. In der Schwelbrandphase lässt sich ein Feuer oft noch mit Wasser, Feuerlöschern oder durch Ersticken löschen. Bleit ein solcher Entstehungsbrand jedoch unbehandelt, kann er sich sehr schnell entwickeln. Sobald das Feuer Nahrung findet, beispielsweise durch andere trockene Materialien, steigt die Temperatur rasant an, der Rauch wird dichter und es kommt zu einem Kronenfeuer, das sich schnell zu einem Vollbrand über mehrere Bäume entwickeln kann, der sich unkontrolliert ausbreitet. Dürreperioden, trockenen Böden und extreme Abholzung verstärken die Brandgefahr und die Ausbreitung des Feuers.


Ökosysteme, in denen Waldbrände zur natürlichen Entwicklung gehören, sind oft an diese Ereignisse angepasst. Dort spielen Brände eine wichtige Rolle bei der Verjüngung von Wäldern, der Freisetzung von Nährstoffen und der Erhaltung der Biodiversität. Im Westen der USA sorgen die Flammen beispielswese dafür, dass die Zapfen sich aufgrund der Hitze öffnen und ihre Samen verteilen, was zur natürlichen Regeneration der wälder betragt. Die Bäume dieser Wälder sind gut durch ihre Rinde geschützt und können so auch größere Brände überstehen. Allerdings geraten auch hier gefördert durch den Klimawandel immer mehr Brände außer Kontrolle, brennen länger, heftiger und größere Flächen als früher.
Besonders gefährdet sind hingegen Wälder, die durch menschliche Eingriffe entstanden sind, wie z.B. Monokulturen oder Plantagen, da diese oft eine geringere Biodiversität aufweisen und damit weniger widerstandsfähig gegenüber Bränden sind. In Deutschland sind beispielsweise Kiefernwälder aufgrund ihrer ätherischen Öle besonders gefährdert, da diese schneller als naturnahe Wälder austrocknen und leicht brennbares Material bilden.


- Natürlcihe Waldbrände durch Blitzeinschalg, aber ehrs selten

- Waldbrände abhängig von Witterung, Brennmatierial und Tageszeit
- Gefährlich, wenn Waldbrände zu heftig, zu lange, zu großflächig an den falschen Orten und zu ungewöhnlcihen Jahreszeiten

- vor Austrieb neuen Grüns steigt Waldbrandgefahr im Frühjahr an und erreicht Ende April/Anfang Mai erstne höhepunkt, Brandgefahr steigt nochmal in Sommermonaten zwischen Ende Juni bis Ende august
- abhängig von Entstehung, Bodenvegetation, Witterung und waldbeständ knnen Brände in vers. Form auftreten
- Bodenfeur am häufigsten, vollfeuer (auch Kronenfeuer genannt) gefährlichste Arten des Brandes -> benötigen zu entstehn und unterhalt immer energiereiches Bodenfeuer
- Bodenfeuer: Brände, die im Laub- und Streuschicht des Waldbodens brennen und relativ langsam voranschreitne
- Kronen- oder wipfelfeuer hingegen verbreiten sich schnell durch Baumkronen und sind schwer zu kontrollieren
- In De nimmt Zahl der Tage mit hohem Brandrisiko zu
- Ursache für Waldbrände zu 96\% von Menschen asgelöst
Lediglich 4,2\% der Brände werden durch natürlcihe Ursache nwie Blitzeinschlag ausgelöst -> erlsöchen oft schnell wieder, da Gewitter mit Blitz in unseren Breiten meist mit starken Neiderschlägen einhergehen
- 20,6\% zwischen 1992 und 2023 lie0 sich Brandstiftung nachweisen, höherre Anteil 23,8\% durch fahrlässtiges Verhalten
- natürliche Waldökosysteme sind robust, Monukulturen anfällig
- Ökologische Folgelast von waldbränden:
    - Entstehung Vielzahl schädlicher emission von Treibhausgase, Feisntaub und anderen Luftschadstoffen bis hin zu Giften wie Kohlenmonoxid, Furanen und Dioxinen
    - Tier- und Pflanzenort vor Ort wird vernichtet
    - Freigebrannter Boden anfälliger für Erosionen, reicht bereits lkeichtet regen
    - Weggespülte Bodenmaterial, wertvoller, humushalter Oberboden fehlt dem Waldökosysstem und ist potzenill gefährlich für nahe Gewässer
    Erodierter Boden büßt Wasserhaltefähigkeit ein ,was Wiederbewaldung erschwert
    Brandfläche sind wasserundurchlässiger und damit trocjenenr als nicht abgebrannte waldbrandflächen
- Dichte Nadelholzwälder unter 40 Jahre als, insbesondere Kifernwälder sind besonders betroffen -> hier besonders trocken udn durch ätherische Öle der Nadelbäume brennen diese beonsder gut
- Beginnt als Bodenfeuer, gesießst drch Reisig, Äste und trockene Pflanzenreste am Waldboden -> Bodenfeuerr begint Nahrung "hinterher" zu laufen
- Wipfelfeuer entsteht -> Springen Flammen nun von Krone zu Kronen spricht man von Vollfeuer oder totalbrand -> wird idR immer durch Bodenfeuer geschützt -> verliert Bodenfeuer seine Kräfte, fällt Vollfeuer zu Boden und sicht sich dort Nahrung

\chapter{Physikalische Detektionsmerkmale}\label{ch:physikalischedetektionsmerkmale}


Brennbare Materialen werden dabei in Brandklassen eingeteilt, die unterschiedliche Brandarten charakterisieren. Wald- oder Buschbrände gehören dabei zur Brandklasse A. Voraussetzung für die Entzündugn von Holz ist eine thermische Aufbereitung, bei der in der ersten Phase Wasserdampf entweicht. Bei weiterem Temperaturanstieg läuft ein Schwelvorgang ab, bei dem zahlreiche Substantzen entstehen, wie z.B. Kohlenstoffmonoxid, Kohlenstoffdioxid, Methan und Wasserstoff. Diese Substanzen brennen utner gelblicher flammenildung und geben dabei so viel Wärme ab, dass die übriggebliebebne Holzkohle in Brand gesetzt wird. Wenn die gasförmigen Stoffe verbrannt sind, gibt es keine Flammebildung mehr und es bleibt nur noch der Glutbrand der Holzkohle übrig.

\chapter{Anforderungen an Frühwarnsysteme/Definition der Vergleichskriterien}\label{ch:anforderungen}

Echtzeitfähigkeit, Fehlalarmresistenz und Robustheit im Gelände

Detektionslatenz, Fehlalarmrate, Skalierbarkeit und ökonomische Effizienz




